\documentclass{beamer}

\usepackage{amsmath, amsfonts, amssymb}

\usetheme{Antibes}

\usecolortheme{dolphin}

\usefonttheme{professionalfonts}

\setbeamertemplate{navigation symbols}{}

\AtBeginSection[]{
  \begin{frame}
    \vfill
    \centering
    \begin{beamercolorbox}[sep=8pt,center,shadow=true,rounded=true]{title}
      \usebeamerfont{title}\insertsectionhead\par%
    \end{beamercolorbox}
    \vfill
  \end{frame}
}

\title{Math for CS \& AI}
\author{Bobby}
\institute{IIIS, Tsinghua University}
\date{\today}

\begin{document}

\begin{frame}
    \titlepage
\end{frame}

\begin{frame}{Contents}
    \tableofcontents[hideallsubsections]
\end{frame}

\begin{section}{Discrete Probability}

\begin{subsection}{Problem 1}
    \begin{frame}
        \begin{block}{Problem 1}
            Consider \(n\) aliens living on a planet with \(m\) days in a year. Let \(A\) be the event that at least two aliens share the same birthday. Prove that the collision probability $\Pr[A]$ satisfies:
            \[\Pr[A] = 1 - e^{-N/m} + \mathcal{O}\left(\frac{1}{\sqrt{m}}\right)\]

            as \(m \to \infty\), where \(N := n(n-1)/2\).    
        \end{block}
    \end{frame}

    \begin{frame}
        \begin{block}{Hint 1}
            It is much easier to first calculate the probability of no collisions, denoted as $\bar{A}$.
            \[
                \Pr[A] = 1 - \Pr[\bar{A}] = 1 - \prod_{k=0}^{n-1} \left(1 - \frac{k}{m}\right)
            \]
        \end{block}

        \pause
        \vspace{2em}
        \begin{block}{Hint 2}
            The asymptotic behavior of the error term depends heavily on the growth rate of $n$. Decompose the proof into two cases based on a threshold:
            \begin{itemize}
                \item Case 1: $n \le m^{0.6}$ (Small $n$)
                \item Case 2: $n > m^{0.6}$ (Large $n$)
            \end{itemize}
        \end{block}

    \end{frame}

    \begin{frame}
        \begin{block}{Hint 3}
            Take the natural logarithm of $\Pr[\bar{A}]$ and apply the Taylor expansion $\ln(1-x) = -x + \mathcal{O}(x^2)$:
            \[
                \ln \Pr[\bar{A}] = \sum_{k=0}^{n-1} \ln\left(1 - \frac{k}{m}\right) = \underbrace{-\sum_{k=0}^{n-1} \frac{k}{m}}_{\text{Main term } -N/m} + \underbrace{\mathcal{O}\left(\sum_{k=0}^{n-1} \frac{k^2}{m^2}\right)}_{\text{Error term } \mathcal{O}(n^3/m^2)}
            \]
        \end{block}
    \end{frame} 

    \begin{frame}
        \begin{block}{Hint 4}
            Use the standard inequality $1 - x \le e^{-x}$ for all $x \in \mathbb{R}$ to upper bound the product directly:
            \[
                \Pr[\bar{A}] \le \prod_{k=0}^{n-1} e^{-k/m} = e^{-N/m}
            \]
            Then, show that for large $n$, both $\Pr[\bar{A}]$ and $e^{-N/m}$ become negligibly small and are absorbed by the $\mathcal{O}\left(\frac{1}{\sqrt{m}}\right)$ error bound.
        \end{block}
    \end{frame}
\end{subsection}

\begin{subsection}{Problem 2}
    \begin{frame}
        \begin{block}{Problem 2}
            Consider throwing $n$ balls independently and uniformly at random into $n$ bins. Let $M$ be the maximum load (the maximum number of balls in any single bin). Prove that with high probability, the maximum load $M$ satisfies:
            \[
                M = \mathcal{O}\left( \frac{\ln n}{\ln \ln n} \right)
            \]
            Specifically, show that for a sufficiently large constant $c$, $\Pr\left[ M \ge \frac{c \ln n}{\ln \ln n} \right] \le \frac{1}{n}$.
        \end{block}
    \end{frame}

    \begin{frame}
        \begin{block}{Hint 1}
            Focus on a single bin, say Bin $i$. Let $X_i$ be the number of balls in Bin $i$. This follows a Binomial distribution $B(n, 1/n)$. The probability that Bin $i$ contains at least $k$ balls is:
            \[
                \Pr[X_i \ge k] \le \binom{n}{k} \left( \frac{1}{n} \right)^k
            \]
        \end{block}

        \pause
        \vspace{1em}
        \begin{block}{Hint 2}
            Use Stirling's approximation or the simpler inequality $\binom{n}{k} \le \left( \frac{ne}{k} \right)^k$ to simplify the bound:
            \[
                \Pr[X_i \ge k] \le \frac{n^k}{k!} \cdot \frac{1}{n^k} = \frac{1}{k!} \le \left( \frac{e}{k} \right)^k
            \]
        \end{block}
    \end{frame}

    \begin{frame}
        \begin{block}{Hint 3}
            Apply the Union Bound over all $n$ bins to bound the maximum load $M$:
            \[
                \Pr[M \ge k] = \Pr\left[ \bigcup_{i=1}^n \{X_i \ge k\} \right] \le \sum_{i=1}^n \Pr[X_i \ge k] \le n \left( \frac{e}{k} \right)^k
            \]
            We want this probability to be $\le n^{-1}$, which means we need $n (e/k)^k \le n^{-1}$, or equivalently $k (\ln k - 1) \ge 2 \ln n$.
        \end{block}
    \end{frame} 

    \begin{frame}
        \begin{block}{Hint 4}
            To solve $k \ln k \approx \ln n$, let $k = \frac{c \ln n}{\ln \ln n}$. Substitute this into the inequality:
            \[
                \ln k = \ln \left( \frac{c \ln n}{\ln \ln n} \right) = \ln \ln n - \ln \ln \ln n + \ln c
            \]
            Show that for a large enough $c$, the $k \ln k$ term dominates $\ln n$ and satisfies the bound. This result highlights the "power of choices" in randomized algorithms.
        \end{block}
    \end{frame}
\end{subsection}

\begin{subsection}{Problem 3}
    \begin{frame}
        \begin{block}{Problem 3}
            Let $G = (V, E)$ be a graph with $n$ vertices and $m$ edges, where $n \le 2m$. Use the probabilistic method to prove that there exists an independent set in $G$ of size at least:
            \[
                \alpha(G) \ge \frac{n^2}{4m}
            \]
            where $\alpha(G)$ denotes the size of the maximum independent set.
        \end{block}
    \end{frame}

    \begin{frame}
        \begin{block}{Hint 1}
            Construct a random subset $S \subseteq V$ by independently including each vertex $v \in V$ with probability $p$. 
            Let $X$ be the number of vertices in $S$, and $Y$ be the number of edges in the subgraph induced by $S$.
            Show that:
            \[
                \mathbb{E}[X] = np, \quad \mathbb{E}[Y] = mp^2
            \]
        \end{block}

        \pause
        \vspace{1em}
        \begin{block}{Hint 2}
            $S$ is not necessarily an independent set. To obtain an independent set $S'$, for each edge in the induced subgraph, remove one of its endpoints from $S$. 
            The size of the resulting independent set $S'$ satisfies:
            \[
                |S'| \ge X - Y
            \]
            Apply the linearity of expectation to find $\mathbb{E}[|S'|]$.
        \end{block}
    \end{frame}

    \begin{frame}
        \begin{block}{Hint 3}
            We have $\mathbb{E}[|S'|] \ge np - mp^2$. To find the tightest bound, maximize this quadratic function with respect to $p$:
            \[
                f(p) = np - mp^2 \implies f'(p) = n - 2mp = 0
            \]
            Set $p = \frac{n}{2m}$ (Note: $n \le 2m$ ensures $p \le 1$). 
            Substitute $p$ back into the expectation to obtain the final bound $\frac{n^2}{4m}$.
        \end{block}
    \end{frame} 

    \begin{frame}
        \begin{block}{Hint 4}
            By the property of expectation, if the average value of $|S'|$ is $\frac{n^2}{4m}$, then there must exist at least one specific set $S'$ in the sample space such that:
            \[
                |S'| \ge \frac{n^2}{4m}
            \]
            This completes the proof. This method is a powerful example of using \textbf{alteration} to satisfy hard constraints in combinatorics.
        \end{block}
    \end{frame}
\end{subsection}

\end{section}

\begin{section}{Sampling and Concentration}

\begin{subsection}{Problem 1}
    \begin{frame}
        \begin{block}{Problem 1}
            For any $k \ge 3$, prove that there exists a graph $G$ such that its chromatic number $\chi(G) > k$ and its girth $g(G) > k$ (where $g(G)$ is the length of the shortest cycle in $G$).
        \end{block}
    \end{frame}

    \begin{frame}
        \begin{block}{Hint 1}
            Fix $\epsilon < 1/k$. Consider a random graph $G(n, p)$ with $p = n^{\epsilon-1}$. 
            Let $X$ be the number of cycles in $G$ of length at most $k$. 
            Show that the expected number of such short cycles is:
            \[
                \mathbb{E}[X] = \sum_{j=3}^k \frac{(n)_j}{2j} p^j = \mathcal{O}(n^k p^k) = \mathcal{O}(n^{k\epsilon})
            \]
            Since $k\epsilon < 1$, $\mathbb{E}[X] = o(n)$. By Markov's inequality, $\Pr[X \ge n/2] \to 0$ as $n \to \infty$.
        \end{block}
    \end{frame}

    \begin{frame}
        \begin{block}{Hint 2}
            Recall that $\chi(G) \ge n/\alpha(G)$. We need to upper bound the independence number $\alpha(G)$. 
            Show that for $r = \lceil \frac{3}{p} \ln n \rceil$:
            \[
                \Pr[\alpha(G) \ge r] \le \binom{n}{r} (1-p)^{\binom{r}{2}} \le \left[ \frac{ne}{r} e^{-p(r-1)/2} \right]^r
            \]
            Prove that this probability tends to $0$ as $n \to \infty$. Thus, w.h.p., $\alpha(G) < \frac{3}{p} \ln n = 3 n^{1-\epsilon} \ln n$.
        \end{block}
    \end{frame}

    \begin{frame}
        \begin{block}{Hint 3}
            Choose $n$ large enough such that $\Pr[X \ge n/2] < 1/2$ and $\Pr[\alpha(G) \ge n/2k] < 1/2$. 
            There exists a graph $G$ with:
            \begin{itemize}
                \item Fewer than $n/2$ cycles of length $\le k$.
                \item $\alpha(G) < n/2k$.
            \end{itemize}
            To obtain $G'$, remove one vertex from each cycle of length $\le k$. 
            The resulting graph $G'$ has $|V'| \ge n/2$ vertices and \textbf{no cycles of length $\le k$}.
        \end{block}
    \end{frame} 

    \begin{frame}
        \begin{block}{Hint 4}
            For the modified graph $G'$, the girth $g(G') > k$ by construction. 
            The independence number $\alpha(G')$ cannot increase by removing vertices, so $\alpha(G') \le \alpha(G) < n/2k$. 
            Then, the chromatic number of $G'$ satisfies:
            \[
                \chi(G') \ge \frac{|V'|}{\alpha(G')} \ge \frac{n/2}{n/2k} = k
            \]
            With a slightly more careful choice of constants, we ensure $\chi(G') > k$.
        \end{block}
    \end{frame}
\end{subsection}

\begin{subsection}{Problem 2}
    \begin{frame}
        \begin{block}{Problem 2}
            In the Coupon Collector's Problem, let $T$ be the number of trials to collect all $n$ distinct coupons. 
            1. Find the variance $\text{Var}(T)$.
            2. Prove that as $n \to \infty$, the variance satisfies the upper bound:
            \[
                \text{Var}(T) \le \frac{\pi^2}{6} n^2
            \]
        \end{block}
    \end{frame}

    \begin{frame}
        \begin{block}{Hint 1}
            Let $X_i$ be the number of trials to collect the $i$-th new coupon after $i-1$ distinct coupons have been collected. 
            $X_i$ follows a Geometric distribution $G(p_i)$ with success probability:
            \[
                p_i = \frac{n - (i-1)}{n}
            \]
            Since $X_i$ are independent, the total time is $T = \sum_{i=1}^n X_i$, and the variance is:
            \[
                \text{Var}(T) = \sum_{i=1}^n \text{Var}(X_i) = \sum_{i=1}^n \frac{1 - p_i}{p_i^2}
            \]
        \end{block}
    \end{frame}

    \begin{frame}
        \begin{block}{Hint 2}
            Substitute $p_i$ into the variance formula:
            \[
                \text{Var}(T) = \sum_{i=1}^n \frac{1 - \frac{n-i+1}{n}}{\left(\frac{n-i+1}{n}\right)^2} = \sum_{i=1}^n \frac{n(i-1)}{(n-i+1)^2}
            \]
            Let $j = n-i+1$. As $i$ ranges from $1$ to $n$, $j$ ranges from $n$ down to $1$:
            \[
                \text{Var}(T) = \sum_{j=1}^n \frac{n(n-j)}{j^2} = n^2 \sum_{j=1}^n \frac{1}{j^2} - n \sum_{j=1}^n \frac{1}{j}
            \]
        \end{block}
    \end{frame}

    \begin{frame}
        \begin{block}{Hint 3}
            To obtain the upper bound, we drop the negative term and consider the limit of the Basel series:
            \[
                \text{Var}(T) < n^2 \sum_{j=1}^n \frac{1}{j^2} < n^2 \sum_{j=1}^{\infty} \frac{1}{j^2}
            \]
            Using the Euler's solution to the Basel problem, we know:
            \[
                \sum_{j=1}^{\infty} \frac{1}{j^2} = \frac{\pi^2}{6}
            \]
        \end{block}
    \end{frame}
    
    \begin{frame}
        \begin{block}{Hint 4}
            Therefore, we have the tight asymptotic bound:
            \[
                \text{Var}(T) \le \frac{\pi^2}{6} n^2
            \]
            This result shows that while the expected time $\mathbb{E}[T] \approx n \ln n$, the standard deviation $\sigma(T)$ is of order $\Theta(n)$, meaning the distribution is relatively concentrated for large $n$.
        \end{block}
    \end{frame}
\end{subsection}

\end{section}


\end{document}